% =====================================================================
% Modelo ABNT de projeto de pesquisa
% EguaTeX (www.eguatex.com.br). Licenca CC0 1.0: dominio publico.
%
% Segue a ABNT NBR 15287:2011 e as normas que ela chama: NBR 6023:2025
% (referencias), NBR 6024:2012 (secoes), NBR 6027:2012 (sumario) e
% NBR 10520:2023 (citacoes).
%
% Como usar: preencha os dados abaixo, troque os textos de exemplo pelo
% seu conteudo, coloque as suas obras no arquivo referencias.bib e compile
% com pdfLaTeX e BibTeX (o EguaTeX faz isso sozinho).
% A formatacao fica no arquivo abnt-eguatex.sty.
% =====================================================================
\documentclass[12pt,a4paper,oneside]{article}
\usepackage[projeto]{abnt-eguatex}

% ---------------------------------------------------------------------
% DADOS DO PROJETO
% ---------------------------------------------------------------------
\instituicao{Nome da entidade}     % na capa, quando a entidade pedir
\autor{Nome do autor}
\titulo{Título do projeto}
\subtitulo{subtítulo, se houver}
\local{Cidade}
\ano{2026}
% Tipo de projeto e entidade a que ele e submetido (4.2.1.1 e)
\natureza{Projeto de pesquisa apresentado ao Curso de Nome do Curso da Nome
  da Instituição, como requisito parcial da disciplina Nome da Disciplina.}
\orientador{Orientador(a): Prof.(ª) Dr.(ª) Nome do orientador}

\begin{document}

% Capa (opcional, 4.1.1). Apague a linha abaixo se nao quiser capa.
\imprimircapa

% Folha de rosto (obrigatoria, 4.2.1.1)
\imprimirfolhaderosto

% Listas (opcionais, 4.2.1.2 a 4.2.1.5). Apague as que nao usar.
\listadequadros

\begin{listadesiglas}
  \sigla{ABNT}{Associação Brasileira de Normas Técnicas}
  \sigla{IBGE}{Instituto Brasileiro de Geografia e Estatística}
\end{listadesiglas}

% Sumario (obrigatorio, 4.2.1.6)
\sumario

% =====================================================================
% ELEMENTOS TEXTUAIS (4.2.2)
% O projeto apresenta: tema, problema, hipoteses (quando houver),
% objetivos, justificativa, referencial teorico, metodologia, recursos e
% cronograma. Os nomes e a ordem das secoes abaixo sao uma sugestao.
% =====================================================================
\textual

\section{Introdução}

Aqui vai a apresentação do tema do projeto. O texto usa letra tamanho 12 e espaçamento de 1,5 entre as linhas.

\subsection{Problema}

Aqui vai o problema de pesquisa, de preferência em forma de pergunta.

\subsection{Hipóteses}

Aqui vão as hipóteses, quando couberem.

\subsection{Objetivos}

\subsubsection{Objetivo geral}

Aqui vai o objetivo geral, começando por um verbo no infinitivo.

\subsubsection{Objetivos específicos}

Os objetivos específicos podem ser escritos em alíneas, como mostra o exemplo:
\begin{alineas}
  \item primeiro objetivo específico;
  \item segundo objetivo específico;
  \item terceiro objetivo específico.
\end{alineas}

\subsection{Justificativa}

Aqui vai a justificativa: por que a pesquisa é importante e para quem.

\section{Referencial teórico}

Aqui vai o referencial teórico que embasa o projeto. As citações seguem o sistema autor-data, como em \cite[p.~15]{livro2024} ou \citeonline[p.~15]{livro2024}. No código, a citação é feita com \texttt{\textbackslash cite\{chave\}}, e a obra entra sozinha na lista de referências. Uma citação indireta não leva aspas \cite{artigo2025,dissertacao2021}.

\begin{citacao}
Aqui vai uma citação direta com mais de três linhas. Ela tem recuo de 4 cm da margem esquerda, letra menor que a do texto, espaçamento simples e não leva aspas. A fonte vem no final, entre parênteses, com o sobrenome, o ano e a página \cite[p.~47]{capitulo2023}.
\end{citacao}

\section{Metodologia}

Aqui vai a metodologia: tipo de pesquisa, população ou corpus, instrumentos de coleta e forma de análise dos dados.

\section{Recursos}

Aqui vão os recursos humanos, materiais e financeiros necessários.

\section{Cronograma}

O Quadro~\ref{qua:cronograma} mostra as etapas da pesquisa ao longo dos meses.

\begin{quadro}[htbp]
  \begin{ilustracao}{\textwidth}
    \caption{Cronograma da pesquisa}\label{qua:cronograma}
    \setstretch{1}\centering
    \begin{tabular}{|p{0.36\linewidth}|*{6}{c|}}
      \hline
      \textbf{Etapa} & \textbf{Mês 1} & \textbf{Mês 2} & \textbf{Mês 3} & \textbf{Mês 4} & \textbf{Mês 5} & \textbf{Mês 6} \\ \hline
      Revisão de literatura & X & X &   &   &   &   \\ \hline
      Coleta de dados       &   & X & X &   &   &   \\ \hline
      Análise dos dados     &   &   & X & X &   &   \\ \hline
      Redação final         &   &   &   & X & X & X \\ \hline
    \end{tabular}
    \fonte{elaborado pelo autor (2026).}
  \end{ilustracao}
\end{quadro}

% =====================================================================
% ELEMENTOS POS-TEXTUAIS (4.2.3)
% =====================================================================
% Referencias (obrigatorias; NBR 6023:2025). A lista e montada sozinha
% a partir do arquivo referencias.bib: entram so as obras citadas no texto,
% em ordem alfabetica e no formato da norma.
\bibliography{referencias}

% Apendices e anexos (opcionais, 4.2.3.3 e 4.2.3.4)
\apendice{Roteiro de entrevista}
Aqui vai o instrumento de coleta de dados feito por você. Apague se não houver.

\end{document}
